\documentclass{article}
\usepackage{geometry}
\geometry{margin=1.5cm, vmargin={0pt,1cm}}
\setlength{\topmargin}{-1cm}
\setlength{\headheight}{12.5pt}
\setlength{\paperheight}{29.7cm}
\setlength{\textheight}{25.3cm}

\usepackage[UTF8]{ctex}
\usepackage{fancyhdr}
\usepackage{layout}
\usepackage{hyperref}
\usepackage{xcolor}
\hypersetup{colorlinks=true, linkcolor=blue!60!black, urlcolor=blue!60!black}

% % % % % % % % % % % % % % % % % % % % % % % % % % % % % % % %
% Common Macros for LaTeX
% % % % % % % % % % % % % % % % % % % % % % % % % % % % % % % %

% Useful packages
\usepackage{amsfonts}
\usepackage{amsmath}
\usepackage{amssymb}
\usepackage{amsthm}
\usepackage{enumerate}
\usepackage{graphicx}
\usepackage{multicol}
\usepackage[ruled,linesnumbered]{algorithm2e}

% Math operators and common symbols
\newcommand{\dif}{\mathrm{d}}
\newcommand{\innerProd}[1]{\left\langle #1 \right\rangle}
\newcommand{\difFrac}[2]{\frac{\dif #1}{\dif #2}}
\newcommand{\pdfFrac}[2]{\frac{\partial #1}{\partial #2}}
\newcommand{\T}{\mathrm{T}}
\newcommand{\norm}[1]{\left\| #1 \right\|}
\newcommand{\abs}[1]{\left| #1 \right|}

% Vector and Matrix shortcuts
\newcommand{\vecTwo}[2]{
  \begin{bmatrix}
    #1 \\
    #2
\end{bmatrix}}
\newcommand{\vecThree}[3]{
  \begin{bmatrix}
    #1 \\
    #2 \\
    #3
\end{bmatrix}}
\newcommand{\vecTwoH}[2]{
  \begin{bmatrix}
    #1 & #2
  \end{bmatrix}
}
\newcommand{\vecThreeH}[3]{
  \begin{bmatrix}
    #1 & #2 & #3
  \end{bmatrix}
}

% Common fields
\newcommand{\R}{\mathbb{R}}
\newcommand{\C}{\mathbb{C}}
\newcommand{\Z}{\mathbb{Z}}
\newcommand{\N}{\mathbb{N}}

% Solutions and proofs
\newenvironment{solution}{
\begin{proof}[解]}{
\end{proof}}

% Utility to get environment variables (requires lualatex)
\newcommand{\getEnv}[2]{\directlua{tex.print(os.getenv("#1") or "#2")}}


% amsthm theorem environments are not used in this document
% (ctex pre-defines \theorem, \definition etc.)

\newcommand{\starKey}{\textcolor{red!65!black}{\textbf{$\star$}}}
\newcommand{\examKey}{\textcolor{red!65!black}{\textbf{[High-frequency exam topic]}}}

\begin{document}

\pagestyle{fancy}
\fancyhead{}
\lhead{Point Set Topology --- Key Concepts \& Theorems}
\rhead{\today}

\begin{center}
  {\Large\bfseries Point Set Topology --- Key Concepts \& Theorems}\\[6pt]
  {\large 点集拓扑重要概念与定理汇总}
\end{center}
\vspace{6pt}

\noindent Compiled from \texttt{report.tex} (Review Outline \& Reference) and
\texttt{all-homework-en.tex} (All Homeworks).
\starKey{} marks high-priority exam topics; \examKey{} marks especially frequent exam topics.
Brief Chinese translations are provided for key concepts and theorems.

\tableofcontents
\newpage

%=====================================================================
\section{Set Theory \& Countability}
%=====================================================================

\subsection{Key Concepts}

\begin{itemize}
  \item \textbf{ZFC Axioms} (ZFC 公理系统): Extensionality, Empty Set, Pairing,
  Union, Power Set, Separation, Replacement, Infinity, Regularity, Choice (AC).
  \item \textbf{Russell's Paradox} (罗素悖论): Unrestricted comprehension
  $\{x\mid P(x)\}$ leads to contradiction; ZFC avoids this via the Separation axiom.
  \item \textbf{Countable} (可数): Finite or countably infinite;
  $\N$, $\Z$, $\Q$, $\N\times\N$ are countable.
  \item \textbf{Uncountable} (不可数): $\R$, $\{0,1\}^\N$, $\mathcal{P}(\N)$ are uncountable.
  \item \textbf{Cardinality} (基数): $\aleph_0=|\N|$, $\mathfrak{c}=|\R|=|\mathcal{P}(\N)|$.
  \item \textbf{Continuum Hypothesis --- CH} (连续统假设): No set $A$ satisfies
  $\aleph_0<|A|<\mathfrak{c}$; independent of ZFC (G\"odel 1940, Cohen 1963).
  \item \textbf{Algebraic numbers} (代数数): countable;
  \textbf{Transcendental numbers} (超越数): uncountable (HW2 Q1).
\end{itemize}

\subsection{Key Theorems \& Proof Techniques}

\begin{itemize}
  \item \starKey \textbf{Cantor's diagonal argument} (Cantor 对角线论证):
  $\{0,1\}^\N$ is uncountable.
  Assume a countable list $x_1,x_2,\ldots$, construct $y$ with $y_i\neq x_{ii}$;
  then $y$ is not in the list.
  \item \textbf{Cantor's Theorem} (Cantor 定理): $|A|<|\mathcal{P}(A)|$ for any set $A$.
  Proof: suppose $g:A\to\mathcal{P}(A)$ is surjective, consider
  $B=\{a\in A:a\notin g(a)\}$, derive a contradiction.
  \item Strategy for countability: construct an injection into $\N$; or write as
  a countable union of countable (finite) sets.
  \item Strategy for uncountability: diagonalization; or inject from a known uncountable set.
\end{itemize}

%=====================================================================
\section{Topological Spaces \& Bases}
%=====================================================================

\subsection{Key Concepts}

\begin{itemize}
  \item \textbf{Topology} (拓扑) $\mathcal{T}\subseteq\mathcal{P}(X)$: contains
  $\varnothing, X$; closed under arbitrary unions; closed under finite intersections.
  \item \textbf{Open set} (开集): $U\in\mathcal{T}$;
  \textbf{Closed set} (闭集): complement is open;
  \textbf{Clopen set} (既开又闭集): both open and closed.
  \item \starKey ``Open/closed'' is always relative to the topology --- never absolute!
  \item \textbf{Discrete topology} (离散拓扑): $\mathcal{T}=\mathcal{P}(X)$.
  \textbf{Indiscrete (trivial) topology} (平凡拓扑/密着拓扑): $\mathcal{T}=\{\varnothing,X\}$.
  \item \textbf{Cofinite topology} (余有限拓扑):
  $\mathcal{T}=\{\varnothing\}\cup\{X\setminus F:F\text{ finite}\}$.
  \item \textbf{Finer / Coarser} (更细/更粗): $\mathcal{T}_1\subseteq\mathcal{T}_2$
  means $\mathcal{T}_2$ is finer (more open sets).
  \item \textbf{Basis} (基) $\mathcal{B}$: (1) covers $X$;
  (2) if $x\in B_1\cap B_2$ then $\exists B_3\ni x$ with $B_3\subseteq B_1\cap B_2$.
  Open sets = arbitrary unions of basis elements.
  \item \textbf{Subbasis} (子基) $\mathcal{S}$: finite intersections form a basis.
  \item \textbf{First-countable} (第一可数): each point has a countable neighborhood basis;
  sequences suffice for topology.
  \item \textbf{Second-countable} (第二可数): the whole space has a countable basis;
  $\R$ has standard basis $\{(a,b):a,b\in\Q\}$ of cardinality $\aleph_0$.
\end{itemize}

\subsection{Key Theorems}

\begin{itemize}
  \item \starKey Second-countable $\Rightarrow$ first-countable. Converse false:
  an uncountable set with discrete topology is first-countable but not second-countable.
  \item Metrizable $\Rightarrow$ first-countable ($\{B(x,1/n)\}$ is a countable local basis).
  \item \textbf{Intersection of topologies is a topology}; union of topologies need not be
  (take it as a subbasis to generate the smallest containing topology).
  \item Topologies on a finite set $\leftrightarrow$ preorders (HW3 Q3).
  \item $\R_\ell$ (lower limit topology, basis $\{[a,b)\}$): first-countable, separable,
  but \textbf{not} second-countable; $\R_\ell\times\R_\ell$ is not normal.
\end{itemize}

%=====================================================================
\section{Metric Spaces}
%=====================================================================

\subsection{Key Concepts}

\begin{itemize}
  \item \textbf{Metric} (度量) $d$: non-negativity, symmetry, triangle inequality,
  $d(x,y)=0\iff x=y$.
  \item \textbf{Pseudometric} (伪度量): may have $d(x,y)=0$ even if $x\neq y$.
  \item \textbf{Open ball / Closed ball} (开球/闭球): $B(x,\varepsilon)$, $\bar{B}(x,\varepsilon)$.
  \item \textbf{Bounded set} (有界集): $\operatorname{diam}(A)<\infty$.
  \item \textbf{Bounded metric} (有界度量): $d^*(x,y)=\min\{d(x,y),1\}$
  induces the same topology (HW5 Q1).
  \item \textbf{Equivalent metrics} (等价度量): induce the same topology;
  sufficient condition: $\exists c_1,c_2>0$ such that $c_1d_1\le d_2\le c_2d_1$.
  \item \textbf{Metrizable} (可度量化): there exists a metric $d$ with $\mathcal{T}=\mathcal{T}_d$.
  \item Three classical metrics on $\R^n$: Euclidean ($\ell^2$), Taxicab ($\ell^1$),
  Supremum ($\ell^\infty$); all equivalent in finite dimensions.
\end{itemize}

\subsection{Key Theorems}

\begin{itemize}
  \item \starKey Metrizable $\Rightarrow$ $T_4$ (normal), $T_2$ (Hausdorff),
  first-countable (these are necessary conditions --- use them to disprove metrizability).
  \item \textbf{Urysohn Metrization Theorem} (Urysohn 可度量化定理):
  $T_3$ + second-countable $\Rightarrow$ metrizable (embed into the Hilbert cube $[0,1]^\omega$).
  \item Compact + metrizable $\Rightarrow$ second-countable
  (compact $\Rightarrow$ totally bounded $\Rightarrow$ separable).
  \item \textbf{For compact spaces: metrizable $\iff$ second-countable $\iff$ separable}.
  \item A pseudometric space is $T_0$ $\iff$ it is a metric space (HW9 Q3).
  \item \examKey Quick reference --- metrizability of common spaces:
  \begin{center}
    \begin{tabular}{@{}lll@{}}
      \hline
      \textbf{Space} & \textbf{Metrizable?} & \textbf{Reason} \\
      \hline
      $\R^n$, standard topology & Yes & $d_2(x,y)=\|x-y\|$ \\
      Discrete topology on any set & Yes & Discrete metric \\
      Indiscrete topology ($|X|\ge2$) & No & Not Hausdorff \\
      Cofinite topology on infinite set & No & Not Hausdorff \\
      $\R^\omega$, product topology & Yes & $d_\omega=\sup_i d^*(x_i,y_i)/(i+1)$ \\
      $\R^\omega$, box topology & No & Not first-countable \\
      \hline
    \end{tabular}
  \end{center}
\end{itemize}

%=====================================================================
\section{Subspace, Product \& Box Topologies}
%=====================================================================

\subsection{Key Concepts}

\begin{itemize}
  \item \textbf{Subspace topology} (子空间拓扑): $\mathcal{T}_Y=\{Y\cap U:U\in\mathcal{T}\}$.
  Open in $Y$ $\neq$ open in $X$ --- always specify the ambient space!
  \item \textbf{Hereditary property} (遗传性质): $X$ has it $\Rightarrow$ every subspace has it.
  Examples: $T_0,T_1,T_2$, second-countable.
  Compactness is \textbf{not} hereditary (only to closed subspaces).
  \item \starKey \textbf{Product topology} (积拓扑): basis $\prod U_i$,
  where $U_i=X_i$ for all but finitely many $i$ (only finitely many coordinates restricted).
  \item \textbf{Box topology} (箱拓扑): basis $\prod U_i$,
  each $U_i$ open --- every coordinate may be restricted.
  \item \textbf{Projection map} (投影映射) $\pi_i$: continuous and open in both topologies.
  \item \textbf{Universal property} (泛性质) --- product topology only:
  $f:Y\to\prod X_i$ is continuous $\iff$ each coordinate $\pi_i\circ f$ is continuous.
  \item \textbf{Eventually-zero sequences} (终零序列) $\R^\omega_{\mathrm{ez}}$:
  sequences with only finitely many nonzero terms.
  \item \textbf{Bounded sequences} (有界序列): $\{x:\exists C,\forall i,|x_i|\le C\}$.
\end{itemize}

\subsection{Key Comparisons \& Theorems}

\begin{itemize}
  \item \starKey For \textbf{finite} products: product topology = box topology.
  For \textbf{infinite} products: box topology is \textbf{strictly finer}.
  \item \examKey Product vs.\ box topology on $\R^\omega$ --- critical comparison:
  \begin{center}
    \begin{tabular}{@{}lcc@{}}
      \hline
      \textbf{Property} & \textbf{Product $\mathcal{T}_p$} & \textbf{Box $\mathcal{T}_b$} \\
      \hline
      Metrizable & Yes ($d_\omega$) & No (not first-countable) \\
      Connected & Yes (path connected) & No (bounded/unbounded separation) \\
      $\operatorname{int}(\R^\omega_{\mathrm{ez}})$ & $\varnothing$ & $\varnothing$ \\
      $\overline{\R^\omega_{\mathrm{ez}}}$ & $\R^\omega$ (dense) & $\R^\omega_{\mathrm{ez}}$ (closed) \\
      Universal property & Holds & Fails \\
      \hline
    \end{tabular}
  \end{center}

  \item Product of connected spaces is connected (in product topology).
  \item Countable product of separable spaces is separable (HW9 Q1).
  \item In box topology, bounded and unbounded sequences form a nontrivial separation,
  so $(\R^\omega,\mathcal{T}_b)$ is disconnected (HW5 Q4).
  \item \textbf{Box topology is not first-countable}: diagonal argument ---
  suppose $\{U_n\}$ is a countable neighborhood basis of $\mathbf{0}$, each
  $U_n$ contains $\prod_i(-\varepsilon_{n,i},\varepsilon_{n,i})$;
  construct $V=\prod_i(-\varepsilon_{i,i}/2,\varepsilon_{i,i}/2)$;
  then no $U_n\subseteq V$.
\end{itemize}

%=====================================================================
\section{Closure, Interior, Boundary \& Limit Points}
%=====================================================================

\subsection{Key Concepts}

\begin{itemize}
  \item \textbf{Closure} (闭包) $\overline{A}$: smallest closed superset of $A$;
  $x\in\overline{A}\iff$ every neighborhood of $x$ meets $A$.
  \item \textbf{Interior} (内部) $A^\circ$: largest open subset of $A$.
  \item \textbf{Boundary} (边界) $\partial A=\overline{A}\setminus A^\circ$;
  $x\in\partial A\iff$ every neighborhood meets both $A$ and $X\setminus A$.
  \item \textbf{Neighborhood} (邻域): a set containing an open set that contains $x$.
  \item \textbf{Limit point / Accumulation point} (极限点/聚点):
  every neighborhood of $x$ meets $A\setminus\{x\}$.
  \item \textbf{Isolated point} (孤立点): $x\in A$ and some neighborhood contains
  no other point of $A$.
  \item \textbf{Derived set} (导集) $A'$ or $LP(A)$: the set of all limit points of $A$.
  \item \textbf{Dense set} (稠密集): $\overline{A}=X$
  (equivalently, meets every nonempty open set).
  \item \textbf{Nowhere dense} (疏集/无处稠密): $\operatorname{int}(\overline{A})=\varnothing$.
  Note: $\Q$ is \textbf{not} nowhere dense in $\R$ ($\overline{\Q}=\R$ has nonempty interior).
  \item \textbf{Regularly open} (正则开集): $A=(\overline{A})^\circ$.
\end{itemize}

\subsection{Key Theorems}

\begin{itemize}
  \item \starKey $\overline{A}=A\cup LP(A)$ (closure = set $\cup$ limit points) --- HW4 Q5.
  \item $\overline{A\cup B}=\overline{A}\cup\overline{B}$.
  \item $A\subseteq B\Rightarrow\overline{A}\subseteq\overline{B}$;
  $\overline{\overline{A}}=\overline{A}$ (idempotence).
  \item $\operatorname{int}(A\cap B)=\operatorname{int}(A)\cap\operatorname{int}(B)$.
  \item \textbf{Kuratowski closure axioms} (Kuratowski 闭包公理):
  If $f:\mathcal{P}(X)\to\mathcal{P}(X)$ satisfies
  $f(\varnothing)=\varnothing$, $A\subseteq f(A)$, $f(f(A))=f(A)$,
  $f(A\cup B)=f(A)\cup f(B)$, then $f(A)=\overline{A}^{\mathcal{T}_f}$
  defines a topology (HW4 Q2).
  \item $(\overline{A})^\circ$ is always regularly open (HW4 Q4).
  \item In a first-countable space, $x\in\overline{A}\iff\exists$ a sequence in $A$
  converging to $x$.
\end{itemize}

%=====================================================================
\section{Continuous Maps \& Homeomorphisms}
%=====================================================================

\subsection{Key Concepts}

\begin{itemize}
  \item \starKey \textbf{Continuity} (连续性) --- four equivalent definitions:
  \begin{enumerate}
    \item $f^{-1}(\text{open})$ is open (global continuity).
    \item Pointwise: $\forall V\ni f(x),\exists U\ni x: f(U)\subseteq V$.
    \item $f^{-1}(\text{closed})$ is closed.
    \item $f(\overline{A})\subseteq\overline{f(A)}$ for all $A$.
  \end{enumerate}
  In metric spaces, (2) reduces to the $\varepsilon$-$\delta$ formulation.
  \item \textbf{Homeomorphism} (同胚): continuous bijection with continuous inverse.
  Denoted $X\cong Y$.
  \item \textbf{Topological invariant} (拓扑不变量): property preserved by homeomorphism
  (compactness, connectedness, Hausdorff, cardinality, etc.).
  \item \textbf{Embedding} (嵌入): $f:X\to Y$ is a homeomorphism onto its image
  (with subspace topology).
  \item \textbf{Pasting Lemma} (粘接引理): $X=A\cup B$ with $A,B$ both closed (or both open),
  $f$ continuous on each piece and agrees on $A\cap B$ $\Rightarrow$ $f$ continuous on $X$.
  \item \textbf{Open map / Closed map} (开映射/闭映射): $f$ sends open/closed sets
  to open/closed sets.
\end{itemize}

\subsection{Six Techniques for Proving/Disproving Homeomorphism}

\begin{enumerate}
  \item \textbf{Explicit construction} (显式构造): write down $f$ and $f^{-1}$,
  verify both are continuous. E.g.\ $(0,1)\cong\R$ via $f(x)=\tan(\pi(x-1/2))$.
  \item \textbf{Topological invariants} (拓扑不变量): find a property $X$ has
  but $Y$ lacks. E.g.\ $[0,1]$ compact, $(0,1)$ not.
  \item \starKey \textbf{Point-removal argument} (去点论证):
  If $X\cong Y$, then $X\setminus\{p\}\cong Y\setminus\{f(p)\}$.
  E.g.\ $\R\not\cong\R^2$ (removing a point disconnects one but not the other).
  \item \textbf{Open/closed mapping criterion} (开/闭映射准则):
  $f$ is a homeomorphism $\iff$ $f$ is a continuous bijection and is open (or closed).
  \item \starKey \textbf{Compact $\to$ Hausdorff shortcut} (紧致到Hausdorff的捷径):
  A continuous bijection from a compact space to a Hausdorff space is automatically
  a homeomorphism --- the most-used tool!
  Proof: $F\subseteq X$ closed $\Rightarrow$ $F$ compact $\Rightarrow$
  $f(F)$ compact $\Rightarrow$ $f(F)$ closed, so $f$ is a closed map.
  \item \textbf{Show a continuous bijection is NOT a homeomorphism}: find an open set
  whose image is not open. E.g.\ $f:[0,1)\to S^1$, $f(t)=(\cos2\pi t,\sin2\pi t)$:
  $[0,1/2)$ is open but its image is not open in $S^1$.
\end{enumerate}

\subsection{Key Theorems}

\begin{itemize}
  \item \examKey Compact + Hausdorff + continuous bijection $\Rightarrow$ homeomorphism
  (one of the most frequently tested facts).
  \item Composition of continuous maps is continuous; restriction of a continuous map
  is continuous.
  \item $f:Y\to\prod X_i$ (product topology) is continuous $\iff$
  each $\pi_i\circ f$ is continuous (universal property).
\end{itemize}

%=====================================================================
\section{Connectedness \& Path Connectedness}
%=====================================================================

\subsection{Key Concepts}

\begin{itemize}
  \item \textbf{Connected} (连通): no separation into two nonempty disjoint open sets;
  equivalently, the only clopen subsets are $\varnothing$ and $X$.
  \item \textbf{Disconnected} (不连通): there exists a \textbf{separation} (分离).
  \item \textbf{Totally disconnected} (完全不连通): only connected subsets are
  singletons (e.g.\ $\Q$, Cantor set).
  \item \textbf{Connected component} (连通分支): maximal connected subset; always closed.
  \item \textbf{Path connected} (道路连通): $\forall a,b\in X$, $\exists$ continuous
  $\gamma:[0,1]\to X$ with $\gamma(0)=a$, $\gamma(1)=b$.
  \item \textbf{Path component} (道路连通分支): equivalence class under
  ``can be joined by a path''; contained in a connected component.
  \item \textbf{Locally path connected} (局部道路连通): every point has a
  path-connected neighborhood basis.
  \item \starKey \textbf{Topologist's sine curve} (拓扑学家的正弦曲线):
  $\{(x,\sin(1/x)):x\in(0,1]\}\cup(\{0\}\times[-1,1])$
  --- connected but NOT path connected. The canonical counterexample.
\end{itemize}

\subsection{Key Theorems}

\begin{itemize}
  \item \starKey \textbf{Path connected $\Rightarrow$ connected}.
  Converse false (see topologist's sine curve).
  \item \starKey \textbf{Connected + locally path connected $\Rightarrow$ path connected}
  (HW6 Q4).
  Proof: Fix $x$, let $S_x=\{y:\exists\text{ path from }x\text{ to }y\}$.
  Show $S_x$ is clopen and nonempty, so $S_x=X$ by connectedness.
  \item \textbf{Continuous image of connected is connected}.
  \item \textbf{Closure of connected is connected}.
  \item \textbf{Union of connected sets with nonempty intersection is connected}.
  \item \textbf{Product of connected spaces is connected} (in product topology).
  \item \textbf{Connected subsets of $\R$ are exactly intervals}.
  Proof: If $A$ is not an interval, $\exists a<c<b$ with $a,b\in A$, $c\notin A$;
  then $A\cap(-\infty,c)$ and $A\cap(c,\infty)$ form a separation.
  \item $\R^n$ is connected; $\R^n\setminus\{0\}$ is connected for $n\ge2$ (path connected).
  \item \starKey \textbf{A connected metric space with more than one point is uncountable}
  (HW7 Q2).
  Proof: pick distinct $a,b$, consider $d_a(x)=d(a,x)$; its continuous image is connected
  in $\R$, hence contains the interval $[0,d(a,b)]$, so it is uncountable.
  \item \textbf{Every continuous $f:[0,1]\to[0,1]$ has a fixed point}
  (1D Brouwer; HW6 Q3).
  Proof: define $g(x)=f(x)-x$, use IVT.
  \item Techniques to prove disconnectedness:
  \begin{enumerate}
    \item Construct a separation $U,V$ directly.
    \item Find a continuous surjection $f:X\to\{0,1\}$ (discrete);
    then $f^{-1}(0),f^{-1}(1)$ form a separation.
    \item Find continuous $f:X\to Y$ with $f(X)$ disconnected
    (contrapositive of ``continuous image of connected is connected'').
  \end{enumerate}
  \item $(0,1)$, $(0,1]$, $[0,1]$ are pairwise non-homeomorphic:
  they have different numbers of non-cut points (0, 1, 2 respectively) --- HW6 Q2.
\end{itemize}

%=====================================================================
\section{Compactness}
%=====================================================================

\subsection{Key Concepts}

\begin{itemize}
  \item \starKey \textbf{Compact} (紧致): every open cover has a finite subcover.
  \item \textbf{Open cover} (开覆盖) / \textbf{Finite subcover} (有限子覆盖):
  core definitions for all compactness proofs.
  \item \textbf{Sequentially compact} (列紧): every sequence has a convergent subsequence
  (equivalent to compact in metric spaces).
  \item \textbf{Locally compact} (局部紧致): every point has a compact neighborhood.
  \item \textbf{One-point compactification} (一点紧化): add $\infty$ to a locally compact
  Hausdorff space to make it compact.
  \item \textbf{Heine--Borel Theorem} (Heine--Borel 定理):
  $A\subseteq\R^n$ is compact $\iff$ closed and bounded.
\end{itemize}

\subsection{Key Theorems --- Highest Exam Priority}

\begin{itemize}
  \item \starKey \textbf{Closed subset of compact is compact}.
  Proof: given an open cover of $F$, add $X\setminus F$ to cover $X$; use compactness of $X$.

  \item \starKey \textbf{Continuous image of compact is compact}.
  Proof: pull back an open cover of $f(X)$ to $X$, extract finite subcover, push forward.

  \item \starKey \textbf{Compact subset of a Hausdorff space is closed}.
  Proof: fix $y\notin K$, for each $x\in K$ use Hausdorff to get disjoint
  $U_x\ni x$, $V_x\ni y$; $\{U_x\}$ covers $K$, take finite subcover;
  $V=\bigcap V_{x_i}$ is an open neighborhood of $y$ disjoint from $K$.

  \item \starKey \textbf{Continuous bijection from compact to Hausdorff is a homeomorphism}
  (see \S6 --- the single most-used tool in point-set topology).

  \item \textbf{Tube Lemma} (管引理): $Y$ compact, if open $N$ contains the slice
  $\{x_0\}\times Y$, then $\exists$ open $W\ni x_0$ with $W\times Y\subseteq N$.
  Proof: for each $y$, pick $U_y\times V_y\subseteq N$; $\{V_y\}$ covers $Y$,
  take finite subcover; set $W=\bigcap U_{y_i}$.

  \item \textbf{Disjoint compact sets in a Hausdorff space can be separated by
  disjoint open sets} (HW7 Q1).

  \item \textbf{Compact + Hausdorff $\Rightarrow$ normal ($T_4$)}.

  \item \textbf{Compact + Hausdorff $\Rightarrow$ regular ($T_3$)}.

  \item Finite product of compact spaces is compact (Tychonoff, finite case).

  \item Arbitrary product of compact spaces is compact (Tychonoff theorem, needs AC).

  \item \examKey \textbf{Every subset is compact in the cofinite topology} (HW6 Q5):
  pick any nonempty $U_0$ from the cover; its complement is finite, so only finitely many
  more sets are needed.

  \item Countable product of finite discrete spaces is compact
  (HW8 Q2: metrizable + diagonal argument for sequential compactness).

  \item $\Z_p$ ($p$-adic integers) is compact (HW8 Q3: closed subset of a product
  of finite discrete spaces).

  \item $\R^\omega$ (product topology) is \textbf{not} locally compact,
  hence has no (Hausdorff) one-point compactification (HW8 Q5).

  \item $\{0\}\cup\{1/n:n\in\N^+\}\cong$ one-point compactification of $\N^+$ (HW8 Q4).

  \item $\Q$ is not locally compact (HW8 Q1).

  \item Equivalent formulation: \textbf{Finite Intersection Property (FIP)} (有限交性质)
  --- in a compact space, any family of closed sets with the FIP has nonempty total intersection.
\end{itemize}

%=====================================================================
\section{Separation Axioms}
%=====================================================================

\subsection{The Hierarchy (层级关系)}

\begin{center}
  $
  T_4\;(\text{normal})\ \Longrightarrow\ T_{3\frac12}\;(\text{Tychonoff})\ \Longrightarrow\
  T_3\;(\text{regular})\ \Longrightarrow\ T_2\;(\text{Hausdorff})\ \Longrightarrow\
  T_1\;(\text{Fr\'echet})\ \Longrightarrow\ T_0\;(\text{Kolmogorov})
  $
\end{center}

Under this reference's convention, $T_3$ and $T_4$ include $T_1$;
none of the reverse implications holds in general.

\subsection{Detailed Breakdown}

\begin{itemize}
  \item \textbf{$T_0$ (Kolmogorov)} (Kolmogorov 空间):
  distinct points are topologically distinguishable (some open set contains exactly one of them).
  Example: Sierpi\'nski space $\{0,1\}$ with topology
  $\{\varnothing,\{1\},\{0,1\}\}$ is $T_0$ but not $T_1$.

  \item \starKey \textbf{$T_1$ (Fr\'echet)} (Fr\'echet 空间):
  singletons are closed (equivalently: two distinct points can be separated one-sidedly).
  Example: cofinite topology on an infinite set is $T_1$ but not $T_2$.
  Properties: hereditary, preserved by products.

  \item \starKey \textbf{$T_2$ (Hausdorff)} (Hausdorff 空间):
  distinct points have disjoint open neighborhoods.
  Consequences: (1) limits are unique; (2) compact subsets are closed;
  (3) the diagonal $\Delta$ is closed in $X\times X$.
  Properties: hereditary, preserved by products.
  Example: any metric space, $\R^n$.
  Non-example: cofinite topology on an infinite set.

  \item \textbf{$T_3$ (Regular)} (正则空间):
  $T_1$ + a point and a disjoint closed set can be separated by disjoint open sets.
  Equivalent: each point has a neighborhood basis of \emph{closed} neighborhoods.
  Properties: hereditary, preserved by products; compact + Hausdorff $\Rightarrow$ regular.

  \item \textbf{$T_{3\frac12}$ (Tychonoff / Completely regular)} (完全正则空间):
  $T_1$ + a point and a disjoint closed set can be separated by a continuous function
  $f:X\to[0,1]$.
  Strictly between $T_3$ and $T_4$; $T_4\Rightarrow T_{3\frac12}\Rightarrow T_3$.

  \item \starKey \textbf{$T_4$ (Normal)} (正规空间):
  $T_1$ + disjoint closed sets can be separated by disjoint open sets.
  \textbf{Crucial:} normality is \textbf{NOT hereditary} (subspaces may fail);
  \textbf{NOT preserved by products} (Sorgenfrey plane $\R_\ell\times\R_\ell$ is not normal).
  Sufficient conditions: compact + Hausdorff $\Rightarrow$ normal; metrizable $\Rightarrow$ normal.
\end{itemize}

\subsection{Core Theorems}

\begin{itemize}
  \item \starKey \textbf{Urysohn's Lemma} (Urysohn 引理):
  $X$ is normal $\iff$ for any disjoint closed $A,B$, there exists continuous
  $f:X\to[0,1]$ with $f(A)=\{0\}$, $f(B)=\{1\}$.
  Proof sketch: construct open sets $U_r$ indexed by dyadic rationals $r$,
  with $\overline{U_r}\subseteq U_s$ for $r<s$; define $f(x)=\inf\{r:x\in U_r\}$.

  \item \textbf{Tietze Extension Theorem} (Tietze 扩张定理):
  $X$ is normal $\iff$ every continuous real-valued function on a closed subset
  extends continuously to $X$.

  \item \textbf{Jones' Lemma} (Jones 引理):
  If a normal space contains a dense set $D$ and a closed discrete set $S$
  with $|S|\ge 2^{|D|}$, then the space is not normal.

  \item \starKey In a $T_3$ space, distinct points can be separated by open sets
  with disjoint closures (HW9 Q4(i)).

  \item In a $T_4$ space, disjoint closed sets can be separated by open sets
  with disjoint closures (HW9 Q4(ii)).

  \item \textbf{Locally compact + Hausdorff $\Rightarrow$ $T_3$ (regular)} --- HW9 Q5.
\end{itemize}

%=====================================================================
\section{Countability Axioms \& Metrizability}
%=====================================================================

\subsection{Key Concepts}

\begin{itemize}
  \item \textbf{Separable} (可分): has a countable dense subset.
  Second-countable $\Rightarrow$ separable; converse false in general (true if metrizable).
  \item \textbf{Lindel\"of}: every open cover has a countable subcover.
  Second-countable $\Rightarrow$ Lindel\"of.
\end{itemize}

\subsection{Key Theorems}

\begin{itemize}
  \item \starKey \textbf{In metric spaces, three properties are equivalent} (HW9 Q2):
  \[
    \text{second-countable}\ \iff\ \text{Lindel\"of}\ \iff\ \text{separable}.
  \]
  Proof outline:
  \begin{itemize}
    \item (i)$\Rightarrow$(ii): countable basis + open cover $\to$
    pick one covering set per basis element.
    \item (ii)$\Rightarrow$(i): for each $n$, $\{B(x,1/n)\}$ is an open cover;
    Lindel\"of gives a countable subcover; union over $n$ yields a countable basis.
    \item (i)$\Rightarrow$(iii): pick one point from each nonempty basis element.
    \item (iii)$\Rightarrow$(i): $\{B(a,1/n):a\in D,\ n\in\N\}$ is a countable basis.
  \end{itemize}

  \item \textbf{Urysohn Metrization Theorem} (Urysohn 可度量化定理):
  $T_3$ + second-countable $\Rightarrow$ metrizable.

  \item For \textbf{compact} spaces: metrizable $\iff$ second-countable $\iff$ separable.

  \item Countable product of separable spaces (product topology) is separable (HW9 Q1).
\end{itemize}

%=====================================================================
\section{Baire Category Theorem}
%=====================================================================

\subsection{Key Concepts}

\begin{itemize}
  \item \textbf{Nowhere dense} (疏集/无处稠密): $\operatorname{int}(\overline{A})=\varnothing$.
  Warning: $\Q$ is \textbf{not} nowhere dense in $\R$ (it's the \emph{closure} that matters:
  $\overline{\Q}=\R$ has nonempty interior).
  \item \textbf{Meager / First category} (第一纲集): countable union of nowhere dense sets.
  \item \textbf{Second category} (第二纲集): not meager.
  \item \textbf{Baire space} (Baire 空间): countable intersection of dense open sets is dense
  (equivalently: countable union of closed nowhere dense sets has empty interior).
\end{itemize}

\subsection{Core Theorems}

\begin{itemize}
  \item \starKey \textbf{Baire Category Theorem} (Baire 纲定理) --- two versions:
  \begin{enumerate}
    \item Every \textbf{complete metric space} is a Baire space.
    \item Every \textbf{locally compact Hausdorff space} is a Baire space.
  \end{enumerate}

  \item \textbf{Special case in compact Hausdorff spaces} (HW7 Q3):
  If $\{A_i\}$ are closed sets with empty interior, then
  $\bigl(\bigcup A_i\bigr)^\circ=\varnothing$.
  Proof: suppose nonempty open $U\subseteq\bigcup A_i$; inductively construct nested
  compact sets avoiding each $A_i$; use FIP to get a contradiction.

  \item \textbf{Cantor set} (Cantor 集) (HW7 Q4):
  compact, totally disconnected, no isolated points, uncountable.

  \item Classical applications of Baire:
  \begin{itemize}
    \item $\R$ is uncountable (otherwise $\R$ would be a countable union of
    nowhere dense singletons, contradicting Baire).
    \item $\Q$ is not a $G_\delta$ set in $\R$.
    \item ``Most'' continuous functions on $[0,1]$ are nowhere differentiable.
  \end{itemize}
\end{itemize}

%=====================================================================
\section{Classic Counterexamples (经典反例)}
%=====================================================================

This section collects the counterexamples that appear repeatedly in the course.
Each one illustrates a specific failure of one property to imply another, or a
boundary case where a theorem's hypotheses cannot be weakened.

\subsection{Connectedness vs.\ Path Connectedness}

\begin{itemize}
  \item \starKey \textbf{Topologist's sine curve} (拓扑学家的正弦曲线):
  \[
    X = \underbrace{\{(x,\sin(1/x)) : x\in(0,1]\}}_{A}\ \cup\
        \underbrace{\{0\}\times[-1,1]}_{B}
  \]
  \begin{itemize}
    \item $A$ is connected (continuous image of $(0,1]$).
    \item $X=\overline{A}$ is connected (closure of connected).
    \item \textbf{Not path connected}: no continuous path can go from a point in $B$
    to a point in $A$, because $\sin(1/x)$ oscillates infinitely as $x\to0^+$.
    \item \textbf{Illustrates}: connected $\nRightarrow$ path connected.
    Also: closure and limit points of a graph (HW4 Q1: $A^\circ=\varnothing$,
    $\overline{A}=A\cup B$, $LP(A)=A\cup B$).
  \end{itemize}
\end{itemize}

\subsection{Separation Axioms}

\begin{itemize}
  \item \textbf{Sierpi\'nski space} ($\{0,1\}$, $\mathcal{T}=\{\varnothing,\{1\},\{0,1\}\}$):
  \textbf{$T_0$ but not $T_1$} (singleton $\{0\}$ is not closed).
  The simplest space distinguishing $T_0$ from $T_1$.

  \item \starKey \textbf{Cofinite topology on an infinite set} (无限集上的余有限拓扑):
  \textbf{$T_1$ but not $T_2$ (Hausdorff)}.
  \begin{itemize}
    \item Every singleton is closed (complement is cofinite), so $T_1$.
    \item Any two nonempty open sets intersect (their complements are finite,
    so the complements' union is finite, hence not the whole space).
    Therefore no two points can be separated by disjoint open sets.
  \end{itemize}

  \item \starKey \textbf{Sorgenfrey line $\R_\ell$} (下限拓扑, basis $\{[a,b)\}$):
  \begin{itemize}
    \item \textbf{$T_3$ (regular)}, first-countable, separable.
    \item \textbf{Not second-countable}: any basis must contain uncountably many sets
    (for each $x$, any basis element containing $x$ must have $x$ as its left endpoint).
    \item \textbf{Illustrates}: separable + first-countable $\nRightarrow$ second-countable.
  \end{itemize}

  \item \starKey \textbf{Sorgenfrey plane $\R_\ell\times\R_\ell$}:
  \textbf{Not $T_4$ (normal)}.
  \begin{itemize}
    \item Each factor $\R_\ell$ is $T_3$ (regular), but the product is not normal.
    \item \textbf{Illustrates}: $T_4$ (normality) is \textbf{not preserved by products};
    $T_3\times T_3$ need not be $T_4$.
  \end{itemize}
\end{itemize}

\subsection{Countability Axioms \& Metrizability}

\begin{itemize}
  \item \textbf{Discrete topology on an uncountable set}:
  \textbf{First-countable} (each $\{x\}$ is a neighborhood basis) but
  \textbf{not second-countable} (would need uncountably many singletons as basis elements).
  \textbf{Illustrates}: first-countable $\nRightarrow$ second-countable.

  \item \textbf{Cocountable topology on an uncountable set} (不可数集上的余可数拓扑):
  $\mathcal{T}=\{\varnothing\}\cup\{X\setminus C:C\text{ countable}\}$.
  \textbf{Not first-countable} (no point has a countable neighborhood basis).
  Hence also not second-countable, not metrizable.

  \item \starKey \textbf{Box topology on $\R^\omega$}:
  \textbf{Not first-countable} (diagonal argument at $\mathbf{0}$), hence
  \textbf{not metrizable}. Also \textbf{disconnected} (bounded/unbounded separation).

  \item \textbf{$\R^\omega$ in the product topology}:
  \textbf{Metrizable} and \textbf{path connected}, but \textbf{not locally compact}.
  \textbf{Illustrates}: metrizable $\nRightarrow$ locally compact.

  \item \textbf{$\R_\ell$ (Sorgenfrey line)}: $T_3$, separable, first-countable, but
  \textbf{not metrizable} (since not second-countable).
  \textbf{Illustrates}: $T_3$ + separable $\nRightarrow$ metrizable;
  the second-countable hypothesis in Urysohn's metrization theorem is essential.

  \item \textbf{Indiscrete (trivial) topology on $|X|\ge2$}:
  \textbf{Not $T_0$}, hence not $T_1$, not $T_2$, \textbf{not metrizable}.

  \item \textbf{Radially open topology on $\R^2$} (径向开集拓扑, HW3 Q5):
  $\mathcal{T}_{\mathrm{std}}\subsetneq\mathcal{T}_{\mathrm{ro}}$.
  \textbf{Strictly finer} than the standard topology; contains sets like a
  ``cross'' through the origin that is not open in the standard topology.
\end{itemize}

\subsection{Compactness}

\begin{itemize}
  \item \starKey \textbf{$(0,1)$, $\R$, $\R^n$}: \textbf{Not compact} (standard topology).
  $[0,1]$ is compact (Heine--Borel); $(0,1)$ is not
  (cover by $(1/n,\,1-1/n)$ has no finite subcover).

  \item \starKey \textbf{Cofinite topology on any set}: \textbf{Every subset is compact}
  (HW6 Q5). In particular, an infinite set with the cofinite topology is compact
  but not Hausdorff. \textbf{Illustrates}: compact $\nRightarrow$ Hausdorff;
  the hypothesis ``Hausdorff'' in ``compact $\Rightarrow$ closed'' is essential.

  \item \textbf{$\R^\omega$ (product topology)}: \textbf{Not locally compact}
  (any basic neighborhood has a free coordinate projecting onto $\R$, which is not compact).
  Hence has no (Hausdorff) one-point compactification (HW8 Q5).

  \item \textbf{$\Q$ with subspace topology}: \textbf{Not locally compact}.
  Any compact subset of $\Q$ has empty interior (otherwise it would contain an interval
  of rationals whose closure in $\R$ contains irrationals).

  \item \textbf{Any infinite set with the discrete topology}:
  Closed and bounded in the discrete metric ($\operatorname{diam}=1$), but
  \textbf{not compact} (cover by singletons has no finite subcover).
  \textbf{Illustrates}: Heine--Borel fails outside $\R^n$; ``closed and bounded''
  depends on the metric.
\end{itemize}

\subsection{Homeomorphisms \& Continuous Maps}

\begin{itemize}
  \item \starKey $f:[0,1)\to S^1$, $f(t)=(\cos2\pi t,\sin2\pi t)$:
  \textbf{Continuous bijection but NOT a homeomorphism} (HW4 Q3).
  $S^1$ is compact, $[0,1)$ is not; or: $f([0,1/2))$ is not open in $S^1$.
  \textbf{Illustrates}: continuous bijection $\nRightarrow$ homeomorphism;
  the compact-to-Hausdorff shortcut is powerful because it handles exactly this gap.

  \item \textbf{$(0,1)\not\cong(0,1]\not\cong[0,1]$} (HW6 Q2):
  Different numbers of non-cut points (0, 1, 2).
  \textbf{Illustrates}: point-removal argument for disproving homeomorphism.

  \item \textbf{$\R\not\cong\R^2$}: $\R\setminus\{p\}$ is disconnected;
  $\R^2\setminus\{p\}$ is connected.
  \textbf{Illustrates}: point-removal + connectedness as a topological invariant.

  \item \textbf{$(0,\infty)\not\cong(0,1]$}: removing any point from $(0,\infty)$
  disconnects it; removing $\{1\}$ from $(0,1]$ leaves $(0,1)$, which is connected.
\end{itemize}

\subsection{Closure, Density \& Product/Box Topology}

\begin{itemize}
  \item \starKey \textbf{Eventually-zero sequences $\R^\omega_{\mathrm{ez}}$} (终零序列):
  \begin{itemize}
    \item In \textbf{product topology}: $\operatorname{int}=\varnothing$,
    $\overline{\R^\omega_{\mathrm{ez}}}=\R^\omega$ (dense).
    \item In \textbf{box topology}: $\operatorname{int}=\varnothing$,
    $\overline{\R^\omega_{\mathrm{ez}}}=\R^\omega_{\mathrm{ez}}$ (closed).
  \end{itemize}
  \textbf{Illustrates}: closure depends on topology; product vs.\ box topology
  give dramatically different closures for the same set (HW5 Q2).

  \item \textbf{Bounded vs.\ unbounded sequences in box topology} (HW5 Q4):
  Both sets are \textbf{open} in $\mathcal{T}_b$, forming a nontrivial separation.
  \textbf{Illustrates}: $(\R^\omega,\mathcal{T}_b)$ is disconnected;
  box topology is ``too fine'' to be connected.

  \item \textbf{Interior of $\Q$ in $\R$}: $\Q^\circ=\varnothing$,
  but $\overline{\Q}=\R$ (dense).
  \textbf{Illustrates}: a set can be dense yet have empty interior;
  $\Q$ is \textbf{not} nowhere dense (the closure has nonempty interior).

  \item \textbf{Nowhere dense vs.\ meager}:
  \begin{itemize}
    \item $\Z$ in $\R$: closed, empty interior $\Rightarrow$ nowhere dense.
    \item $\{1/n:n\in\N^+\}$ in $\R$: closure adds $\{0\}$, still has empty interior
    $\Rightarrow$ nowhere dense.
    \item $\Q$ in $\R$: \textbf{not} nowhere dense ($\overline{\Q}=\R$),
    but \textbf{is} meager ($\Q=\bigcup_{q\in\Q}\{q\}$, countable union of nowhere dense singletons).
  \end{itemize}
  \textbf{Illustrates}: nowhere dense $\neq$ meager;
  meager sets can be dense (and large in that sense).

  \item \textbf{Cantor set $C$} (Cantor 集, HW7 Q4):
  Compact, totally disconnected, no isolated points, uncountable, nowhere dense,
  Lebesgue measure zero.
  \textbf{Illustrates}: a set can be uncountable yet nowhere dense and measure zero;
  a compact metric space can be totally disconnected yet have no isolated points.
\end{itemize}

\subsection{Metrizability --- Necessary Conditions at a Glance}

\begin{center}
  \begin{tabular}{@{}llll@{}}
    \hline
    \textbf{Space} & \textbf{Not metrizable because\ldots} & \textbf{Fails} \\
    \hline
    Indiscrete topology ($|X|\ge2$) & not Hausdorff & $T_2$ \\
    Cofinite topology on infinite set & not Hausdorff & $T_2$ \\
    $\R^\omega$, box topology & not first-countable & first-countable \\
    $\R_\ell$ (Sorgenfrey line) & not second-countable\phantom{X} & (Urysohn needs $2$nd-ctble) \\
    Cocountable topology (uncountable) & not first-countable & first-countable \\
    Sorgenfrey plane & not normal & $T_4$ \\
    \hline
  \end{tabular}
\end{center}

%=====================================================================
\section{Proof Techniques Summary}
%=====================================================================

\begin{enumerate}
  \item \textbf{Double inclusion} (双向包含): prove $A=B$.
  E.g.\ $\overline{A}=A\cup LP(A)$, topology equality $\mathcal{T}_1=\mathcal{T}_2$.

  \item \textbf{Contrapositive} (逆否命题): $P\Rightarrow Q$ proved as $\neg Q\Rightarrow\neg P$.
  E.g.\ continuous image preserves connectedness.

  \item \textbf{Contradiction} (反证法): assume $\neg$ conclusion, derive contradiction.
  E.g.\ Cantor's theorem.

  \item \textbf{Case analysis} (情形分析): split into exhaustive cases.
  E.g.\ triangle inequality for $d^*$ ($a\ge1$ or $b\ge1$ vs.\ both $<1$).

  \item \textbf{Sequence construction} (序列构造): build sequences to show
  $x\in\overline{A}$ or $x\in LP(A)$. E.g.\ closure of topologist's sine curve.

  \item \starKey \textbf{Free coordinates} (自由坐标): in product topology, unrestricted
  coordinates allow counterexample construction.
  E.g.\ $\R^\omega_{\mathrm{ez}}$ has empty interior.

  \item \textbf{Box neighborhood trick} (箱邻域技巧):
  $\prod_i(x_i-\varepsilon_i,x_i+\varepsilon_i)$ controls every coordinate.
  E.g.\ bounded sequences form an open set in box topology.

  \item \textbf{Continuous function to a discrete space}: prove disconnectedness.
  E.g.\ $X=\{(x,0)\}\cup\{(x,1/x)\}$, define $f(x,y)=xy$; $f(X)=\{0,1\}$.

  \item \starKey \textbf{Topological invariant argument} (拓扑不变量论证):
  prove $X\not\cong Y$. E.g.\ $[0,1)\not\cong S^1$ (compactness).

  \item \starKey \textbf{Clopen set argument} (Clopen集论证): in a connected space,
  define $S=\{x:\text{property holds at }x\}$, show $S$ is clopen and nonempty,
  then $S=X$. E.g.\ connected + locally path connected $\Rightarrow$ path connected.

  \item \textbf{IVT / fixed point} (介值定理/不动点): define $g(x)=f(x)-x$,
  check signs at endpoints, apply IVT. E.g.\ 1D Brouwer.

  \item \textbf{Ball comparison} (球比较法): show each $d_1$-ball contains a $d_2$-ball
  and vice versa to prove $\mathcal{T}_{d_1}=\mathcal{T}_{d_2}$.
  E.g.\ $d^*$ induces the same topology as $d$.

  \item \textbf{Monotonicity from additivity} (加法推出单调性):
  $A\subseteq B\Rightarrow f(A)\subseteq f(B)$ from $f(A\cup B)=f(A)\cup f(B)$.
  E.g.\ Kuratowski closure axioms.

  \item \textbf{Characteristic function} (特征函数): encode subsets as functions
  $A\to\{0,1\}$; set operations become pointwise arithmetic.
  E.g.\ $\mathcal{P}(A)\cong 2^A$ (HW1 Q1).

  \item \textbf{Countable union argument} (可数并论证):
  $S=\bigcup_{n=1}^\infty S_n$ with each $S_n$ countable (or finite) $\Rightarrow$
  $S$ countable. E.g.\ algebraic numbers, eventually-zero sequences.

  \item \textbf{Well-definedness} (良定义性): when defining a function on equivalence classes,
  verify independence from the choice of representative.
  E.g.\ quotient metric from pseudometric (HW2 Q3).

  \item \textbf{Axiom verification checklist} (公理验证清单): systematically verify
  each axiom --- topology (3), metric (3), equivalence relation (3), group (4).

  \item \textbf{Subbasis method} (子基方法): the smallest topology/equivalence relation
  containing a family = the intersection of all topologies/relations containing it.

  \item \textbf{Minimal counterexample} (最小反例): ``Can one always\ldots'' questions
  are often resolved by the smallest nontrivial example.
  E.g.\ $\mathcal{T}_1\cup\mathcal{T}_2$ not a topology ($|S|=3$; HW3 Q4).

  \item \textbf{Finite subcover + finiteness} (有限子覆盖 + 有限性): after extracting
  a finite subcover from a compact space, use finiteness (e.g.\ the minimum of finitely
  many positive numbers is still positive).

  \item \textbf{Preimage arguments} (原像论证): $f^{-1}$ commutes with union, intersection,
  and complement (no continuity needed!).
\end{enumerate}

%=====================================================================
\section{Exam Priority}
%=====================================================================

\noindent The following priority order is suggested in \texttt{report.tex} (highest to lowest):

\begin{center}
  \textbf{Compactness $>$ Separation Axioms $>$ Urysohn Lemma / Metrization $>$
  Subspace \& Product Topologies $>$ Connectedness \& Path Connectedness $>$
  Closure / Interior / Continuity $>$ Set Theory \& Countability}
\end{center}

The final exam is described as mainly variations of homework, midterm, handout,
and in-class problems: typically about 5 questions, mostly familiar types plus
possibly one transfer problem.
The target standard is \textbf{being able to write each proof independently},
not merely recognizing a solution.

\end{document}
